\documentclass[tikz,border=8pt]{standalone}
\usepackage{amsmath}
\usepackage{fontspec}
\usetikzlibrary{arrows.meta,calc,positioning}

\setmainfont{Arial}
\setsansfont{Arial}

\definecolor{ink}{HTML}{243447}
\definecolor{muted}{HTML}{607086}
\definecolor{orbit}{HTML}{40566F}
\definecolor{velocity}{HTML}{D1495B}
\definecolor{acceleration}{HTML}{246BCE}
\definecolor{rotation}{HTML}{C47F17}
\definecolor{softblue}{HTML}{EEF5FF}
\definecolor{softred}{HTML}{FFF1F3}
\definecolor{softgold}{HTML}{FFF7E8}
\definecolor{canvas}{HTML}{FFFFFF}

\begin{document}
\begin{tikzpicture}[
  x=1cm,
  y=1cm,
  font=\sffamily\large,
  >={Latex[length=3.2mm,width=2.2mm]},
  vector/.style={-{Latex[length=3.2mm,width=2.2mm]}, line width=1.7pt},
  construction/.style={draw=muted, dashed, line width=0.85pt},
  panel/.style={rounded corners=7pt, line width=0.9pt, minimum width=8.00cm, inner sep=10pt},
  formula/.style={font=\sffamily\Large, text=ink},
  caption/.style={font=\sffamily\normalsize, text=muted, align=left}
]
  \path[fill=canvas] (-0.35,-0.35) rectangle (18.45,10.45);

  \node[anchor=west, text=ink, font=\sffamily\bfseries\LARGE]
    at (0.25,9.95) {Движение материальной точки по окружности};
  \draw[draw=ink!18, line width=0.8pt] (0.25,9.50) -- (17.85,9.50);

  \node[anchor=west, text=muted, font=\sffamily\bfseries\large]
    at (0.35,9.02) {ГЕОМЕТРИЯ ДВИЖЕНИЯ};
  \node[anchor=west, text=muted, font=\sffamily\bfseries\large]
    at (9.85,9.02) {ФИЗИЧЕСКИЙ СМЫСЛ};

  \coordinate (O) at (4.60,4.75);
  \coordinate (A) at ($(O)+(38:3.10)$);
  \coordinate (X) at ($(O)+(3.10,0)$);
  \coordinate (VEnd) at ($(A)+(128:2.35)$);
  \coordinate (AEnd) at ($(A)!0.62!(O)$);

  \fill[softblue!45] (O) circle (3.10);
  \draw[draw=orbit, line width=1.55pt] (O) circle (3.10);

  \draw[construction] (O) -- (X);
  \draw[construction] (O) -- (A);
  \draw[-{Latex[length=2.5mm]}, draw=rotation, line width=1.25pt]
    ($(O)+(0:1.20)$) arc[start angle=0,end angle=38,radius=1.20];
  \node[text=rotation, font=\sffamily\Large] at ($(O)+(19:1.62)$) {$\varphi$};

  \draw[-{Latex[length=3mm]}, draw=rotation, line width=1.65pt]
    ($(O)+(205:3.55)$) arc[start angle=205,end angle=265,radius=3.55];
  \node[text=rotation, font=\sffamily\Large, anchor=east]
    at ($(O)+(230:4.00)$) {$\omega$};

  \draw[vector, draw=velocity] (A) -- (VEnd);
  \node[text=velocity, font=\sffamily\Large, anchor=south east]
    at ($(A)!0.72!(VEnd)+(-0.10,0.10)$) {$\vec v$};

  \draw[vector, draw=acceleration] (A) -- (AEnd);
  \node[text=acceleration, font=\sffamily\Large, anchor=west]
    at ($(A)!0.52!(AEnd)+(0.30,-0.32)$) {$\vec a_{\text{цс}}$};

  \fill[ink] (O) circle (2.7pt);
  \node[text=ink, anchor=north, font=\sffamily\large] at ($(O)+(0,-0.18)$) {$O$};

  \fill[velocity] (A) circle (4pt);
  \node[text=ink, anchor=south west, font=\sffamily\bfseries\large]
    at ($(A)+(0.10,0.12)$) {$A$};

  \node[text=muted, inner sep=1pt, rotate=38]
    at ($(O)!0.40!(A)+(0,0.30)$) {$R$};

  \node[panel, draw=velocity!55, fill=softred, anchor=north west, minimum height=1.48cm]
    (speedpanel) at (9.75,8.55) {};
  \node[anchor=west, text=velocity, font=\sffamily\bfseries\Large]
    at (10.10,7.98) {$\vec v$};
  \node[caption, anchor=west, text width=3.55cm]
    at (10.95,7.98) {направлена по касательной};
  \node[formula, anchor=east] at (17.45,7.98) {$v=\omega R$};

  \node[panel, draw=acceleration!55, fill=softblue, anchor=north west, minimum height=1.58cm]
    (accpanel) at (9.75,6.48) {};
  \node[anchor=west, text=acceleration, font=\sffamily\bfseries\Large]
    at (10.10,5.88) {$\vec a_{\text{цс}}$};
  \node[caption, anchor=west, text width=3.15cm]
    at (11.50,5.88) {направлено к центру};
  \node[formula, anchor=east] at (17.45,5.88)
    {$a_{\text{цс}}=\dfrac{v^2}{R}=\omega^2R$};

  \node[panel, draw=rotation!60, fill=softgold, anchor=north west, minimum height=2.75cm]
    (uniformpanel) at (9.75,4.22) {};
  \node[anchor=west, text=rotation, font=\sffamily\bfseries\large]
    at (10.10,3.53) {РАВНОМЕРНОЕ ДВИЖЕНИЕ};
  \node[anchor=center, text=ink, font=\sffamily\Large]
    at (13.70,2.72) {$|\vec v|=\mathrm{const}$};
  \node[caption, anchor=center, align=center]
    at (13.70,1.92) {модуль скорости постоянен;\\направление вектора меняется};

  \node[anchor=west, text=muted, font=\sffamily\normalsize]
    at (0.45,0.20) {Векторы $\vec v$ и $\vec a_{\text{цс}}$ взаимно перпендикулярны.};
\end{tikzpicture}
\end{document}
